% =============================================================================
% EC4308 Final Project Report
% "Forecasting the Equity Risk Premium Out-of-Sample: Regularized Machine
%  Learning, Technical Indicators & Tactical Asset Allocation"
% Format: <= 15 pages, 11pt, 1.5 line spacing (Prof. Denis's requirements)
% Build:  latexmk -pdf report/main.tex
% =============================================================================
\documentclass[11pt,a4paper]{article}

\usepackage[margin=1in]{geometry}
\usepackage{amsmath,amssymb}
\usepackage{booktabs}
\usepackage{graphicx}
\usepackage{caption}
\usepackage[colorlinks=true,linkcolor=blue,citecolor=blue,urlcolor=blue]{hyperref}
\usepackage{setspace}
\usepackage{threeparttable}
\usepackage{longtable}
\usepackage{pdflscape}

\onehalfspacing

\title{\vspace{-1.2cm}Forecasting the Equity Risk Premium Out-of-Sample:\\[2mm]
Regularized Machine Learning, Technical Indicators and\\Tactical Asset Allocation\vspace{-2mm}}
\author{Leon \quad En Shao \quad Sasinthiran Kannan \quad Vienna\\
EC4308 Machine Learning and Economic Forecasting, AY 2026/2027 Semester 1}
\date{9 November 2026}

\begin{document}
\maketitle
\vspace{-1.6cm}

\begin{abstract}
\noindent We re-examine out-of-sample predictability of the US equity premium using a
fixed 240-month rolling-window design on data spanning February 1951 to December 2025.
The predictor set combines the fourteen Goyal--Welch (2008) macroeconomic variables
(with one to six monthly lags) and the fourteen technical indicators of
Neely, Rapach, Tu and Zhou (2014), which we reproduce from raw price and volume
data with a 99.98\% signal-month match to the published series. Against the
prevailing historical mean benchmark we evaluate $\ell_1$-penalized least
squares (LASSO and post-LASSO), Ridge, Elastic Net, principal components and
partial least squares regressions, and random forests, together with equal-weight
and constrained Granger--Ramanathan forecast combinations, at horizons
$h\in\{1,3,12\}$. Statistical accuracy is assessed with the Diebold--Mariano
(1995) test on squared and absolute errors (Clark--West as an appendix check),
and economic value with a binary switching strategy and a mean--variance
investor, both net of transaction costs. [KEY FINDINGS -- TO BE FILLED AFTER
PHASE 3 EVALUATION.]
\end{abstract}

\section{Introduction}\label{sec:intro}
% WRITTEN LAST per the course guidelines: state research questions, motivation
% (Goyal-Welch 2008/2024 failure of fundamentals; Neely et al. 2014 technical
% indicators; Gu-Kelly-Xiu 2020 ML), summarize results, outline the paper.
[TO BE FILLED AFTER PHASE 3 -- SEE \ref{sec:results}.]

\section{Data}\label{sec:data}

\subsection{Sources and sample}\label{sec:data:sources}
The target variable is the monthly log equity premium on the S\&P 500,
$y_t = 100\,(\ln(1+R^m_t) - \ln(1+R^f_t))$, where $R^m_t$ is the value-weighted
S\&P 500 return including dividends and $R^f_t$ the risk-free rate, taken from
the continuously updated Goyal--Welch--Zafirov (2024) dataset (monthly
observations 1871:01--2025:12). The premium is expressed in percent per month;
over the estimation sample its mean is 0.66\% with a standard deviation of
4.21\%. The fourteen macroeconomic predictors follow Goyal and Welch (2008):
dividend-price ratio, dividend yield, earnings-price ratio, dividend payout
ratio, stock variance, book-to-market, net equity expansion, Treasury bill
rate, long-term yield, long-term return, term spread, default yield spread,
default return spread, and inflation. Each macro predictor enters with one to
six monthly lags.

The fourteen technical indicators follow Neely et al.\ (2014): six moving
average rules $\mathrm{MA}(j,m)$ with $j\in\{1,2,3\}$ and $m\in\{9,12\}$
comparing short and long moving averages of the index level, two momentum
rules $\mathrm{MOM}(m)$, $m\in\{9,12\}$, and six on-balance volume rules
$\mathrm{VOL}(j,m)$ comparing moving averages of on-balance volume. Each
indicator is a binary signal (1 = hold equities). We compute the signals from
the raw index level and monthly S\&P 500 volume: the paper's own volume series
for 1950--2011 and monthly average daily volume from Yahoo Finance for
2012--2025 (the two series have correlation 0.9999 on their overlap, and the
volume-based signals are invariant to a constant volume rescaling). Against
the published signals of Neely et al.\ (2014) over their 1950:12--2011:12
sample, 12 of the 14 signals match exactly and the remaining two match in
99.9\% of months; the two mismatching months are documented data revisions.

\subsection{Alignment and sample period}\label{sec:data:align}
Row $t$ of the panel holds the premium \emph{earned} in month $t$ together
with predictors \emph{observable at the end of month $t$}. Macro lag columns
satisfy $x^{\mathrm{lag}\,k}_t = x_{t-k}$, $k=1,\dots,6$; technical signals at
row $t$ use price and volume data through month $t$ only, so no forecast uses
information unavailable at its origin. The sample is the first month in which
every column is non-missing (1951:02, driven by the twelve-month burn-in of
the volume rules) through 2025:12, giving $T=899$ monthly observations.
Table~\ref{tab:table1} reports summary statistics and augmented
Dickey--Fuller (BIC lag selection, intercept and trend) tests.

\begin{table}[htbp]
\centering
\caption{Summary statistics and unit-root tests (sample: 1951:02--2025:12).}
\label{tab:table1}
\small
\begin{threeparttable}
\begin{tabular}{lrrrrrrrc}
\toprule
Series & Mean & SD & Min & Max & Skew & Kurt & ADF & Sig. \\
\midrule
\multicolumn{9}{l}{\emph{Panel A: log equity premium (\%, monthly)}} \\
$y_t$ (full sample) & 0.662 & 4.211 & -22.18 & 16.30 & -0.42 & 4.54 & -29.44 & 1\% \\
\midrule
\multicolumn{9}{l}{\emph{Panel B: macroeconomic predictors (raw series)}} \\
d/p & 0.030 & 0.013 & 0.01 & 0.07 & 0.69 & 2.91 & -3.71 & 5\% \\
d/y & 0.030 & 0.013 & 0.01 & 0.07 & 0.69 & 2.91 & -3.44 & 10\% \\
e/p & 0.063 & 0.026 & 0.01 & 0.15 & 0.98 & 3.58 & -3.68 & 5\% \\
d/e & 0.500 & 0.283 & 0.29 & 3.97 & 8.48 & 88.54 & -6.71 & 1\% \\
svar & 0.002 & 0.005 & 0.00 & 0.07 & 10.87 & 149.96 & -20.06 & 1\% \\
b/m & 0.486 & 0.250 & 0.12 & 1.21 & 0.70 & 2.72 & -2.06 & >10\% \\
ntis & 0.011 & 0.020 & -0.06 & 0.05 & -0.62 & 2.88 & -4.23 & 1\% \\
tbl & 0.041 & 0.030 & 0.00 & 0.16 & 0.92 & 4.25 & -2.88 & >10\% \\
lty & 0.057 & 0.028 & 0.01 & 0.15 & 0.83 & 3.30 & -1.86 & >10\% \\
ltr & 0.005 & 0.027 & -0.11 & 0.15 & 0.52 & 6.05 & -28.40 & 1\% \\
tms & 0.016 & 0.014 & -0.04 & 0.05 & -0.04 & 2.75 & -4.78 & 1\% \\
dfy & 0.009 & 0.004 & 0.00 & 0.03 & 1.89 & 8.25 & -4.43 & 1\% \\
dfr & 0.000 & 0.014 & -0.10 & 0.07 & -0.66 & 10.97 & -24.01 & 1\% \\
infl & 0.003 & 0.004 & -0.02 & 0.02 & 0.14 & 5.54 & -3.38 & 10\% \\
\midrule
\multicolumn{9}{l}{\emph{Panel C: technical signals (binary; mean = fraction of months in ``buy'')}} \\
MA\_1\_9 & 0.709 & & & & & & & \\
MA\_1\_12 & 0.729 & & & & & & & \\
MA\_2\_9 & 0.710 & & & & & & & \\
MA\_2\_12 & 0.730 & & & & & & & \\
MA\_3\_9 & 0.715 & & & & & & & \\
MA\_3\_12 & 0.731 & & & & & & & \\
MOM\_9 & 0.729 & & & & & & & \\
MOM\_12 & 0.743 & & & & & & & \\
VOL\_1\_9 & 0.694 & & & & & & & \\
VOL\_1\_12 & 0.720 & & & & & & & \\
VOL\_2\_9 & 0.689 & & & & & & & \\
VOL\_2\_12 & 0.719 & & & & & & & \\
VOL\_3\_9 & 0.704 & & & & & & & \\
VOL\_3\_12 & 0.718 & & & & & & & \\
\bottomrule
\end{tabular}
\begin{tablenotes}
\footnotesize
\item ADF: augmented Dickey--Fuller statistic with intercept and trend; lag order by BIC; critical values $-4.04/-3.45/-3.15$ (1\%/5\%/10\%). $T=899$ monthly observations.
\item The premium is $100\,(\ln(1+R^m)-\ln(1+R^f))$; macro predictors enter the models with one to six lags; technical signals are the 14 binary rules of Neely et al. (2014).
\end{tablenotes}
\end{threeparttable}
\end{table}


\section{Methodology}\label{sec:method}

\subsection{Rolling-window forecasting design}\label{sec:method:design}
All models are estimated on a fixed 240-month (20-year) rolling window and
produce an $h$-month-ahead forecast of the premium, $h\in\{1,3,12\}$. The
first out-of-sample forecast targets February 1971, leaving 659 out-of-sample
months at $h=1$. The benchmark is the prevailing (historical) mean premium
computed through the forecast origin, the only admissible benchmark for this
exercise: OLS benchmark forecasts fail out-of-sample and are not used.

\subsection{Models}\label{sec:method:models}
We estimate seven models within each window: (i) LASSO with the theoretical,
data-driven penalty of \texttt{hdm::rlasso} (no cross-validation required);
(ii) post-LASSO, i.e.\ OLS refit on the LASSO-selected support; (iii) Ridge
and (iv) Elastic Net ($\alpha=0.5$) with the penalty chosen by ten-fold
cross-validation \emph{inside the training window} (no information from the
test month is ever used); (v) principal components regression and (vi)
partial least squares with the number of components selected by
cross-validated PRESS (scaled inputs, at most 20 components); and (vii) a
random forest (500 trees, $m_{\mathrm{try}}=p/3$, minimum node size 5, fixed
seed, single-threaded for exact reproducibility). All models use the same
$98$-predictor matrix (84 macro-lag columns plus 14 technical signals); the
expanded Goyal--Welch--Zafirov (2024) predictors are held out as a
robustness variant.

\subsection{Forecast combinations}\label{sec:method:combos}
We combine the seven model forecasts with (i) the simple equal-weight average
and (ii) a constrained Granger--Ramanathan (1984) regression: at each origin
the weights solve $\min_w \|y-Fw\|^2$ subject to $\sum_i w_i = 1$ and
$w_i\ge0$, estimated only on past out-of-sample forecast errors known at that
origin (targets $\le$ origin). Before 60 usable observations accumulate,
equal weights are used.

\subsection{Evaluation}\label{sec:method:eval}
Statistical accuracy is measured by out-of-sample RMSE and MAE and the
Campbell--Thompson (2008) $R^2_{\mathrm{OOS}} = 1 - \sum e_t^2 / \sum
e_{t,\mathrm{PM}}^2$ against the prevailing mean. Significance is tested with
the Diebold--Mariano (1995) statistic on squared and absolute error
differentials, with Newey--West long-run variance (lag $h-1$); the
Clark--West (2007) adjusted-MSPE test is reported in the appendix for the
nested comparison against the benchmark. Economic value is assessed with two
strategies net of transaction costs (baseline 10 bps, sensitivity 5 and
20 bps): (i) a binary switching strategy that holds the S\&P 500 when the
$h$-month-ahead premium forecast is positive and three-month Treasury bills
otherwise, reporting terminal wealth, annualized Sharpe ratio, maximum
drawdown, and turnover; and (ii) a mean--variance investor with risk aversion
$\gamma=3$ who allocates $w_t = (1/\gamma)\,\hat y_t / \hat\sigma^2_t$ to
equities, clamped to $[0,1.5]$, where $\hat\sigma^2_t$ is realized variance
over the previous 60 months; the certainty-equivalent return gain
$\Delta U$ relative to the benchmark allocation is reported.

\section{Results}\label{sec:results}
% Table 2 (statistical accuracy), Figure 2 (R2OOS), Table 3 (economic value),
% Figure 1 (switching wealth), Figure 3 (GR weights) -- filled after Phase 3.
[TO BE FILLED AFTER PHASE 3 EVALUATION.]

\section{Conclusion}\label{sec:conclusion}
[TO BE FILLED AFTER PHASE 3 -- summarize, limitations, future work.]

\appendix
\section{Clark--West appendix test}\label{app:cw}
[TO BE FILLED AFTER PHASE 3.]

\begin{thebibliography}{9}
\bibitem{gw2008} Goyal, A., and I. Welch (2008). A comprehensive look at the
empirical performance of equity premium prediction. \emph{Review of Financial
Studies} 21(4), 1455--1508.
\bibitem{gwz2024} Goyal, A., I. Welch, and A. Zafirov (2024). A comprehensive
2024 look at the empirical performance of equity premium prediction.
\emph{Review of Financial Studies}.
\bibitem{neely2014} Neely, C. J., D. E. Rapach, J. Tu, and G. Zhou (2014).
Forecasting the equity risk premium: The role of technical indicators.
\emph{Management Science} 60(7), 1772--1791.
\bibitem{gkx2020} Gu, S., B. Kelly, and D. Xiu (2020). Empirical asset pricing
via machine learning. \emph{Review of Financial Studies} 33(5), 2223--2273.
\bibitem{ct2008} Campbell, J. Y., and S. B. Thompson (2008). Predicting excess
stock returns out of sample: Can anything beat the historical average?
\emph{Review of Financial Studies} 21(4), 1509--1531.
\bibitem{dm1995} Diebold, F. X., and R. S. Mariano (1995). Comparing predictive
accuracy. \emph{Journal of Business \& Economic Statistics} 13(3), 253--263.
\bibitem{cw2007} Clark, T. E., and K. D. West (2007). Approximately normal
tests for equal predictive accuracy in nested models. \emph{Journal of
Econometrics} 138(1), 291--311.
\bibitem{gr1984} Granger, C. W. J., and R. Ramanathan (1984). Improved methods
of combining forecasts. \emph{Journal of Forecasting} 3(2), 197--204.
\end{thebibliography}

\end{document}
